% =====================================================================
%  CV general en ESPAÑOL. Cuerpo; el preámbulo está en cv-preamble.tex.
%
%  No compilar este archivo directamente. Usar:
%    johan-rodriguez-cv-es-onepage.tex → 1 página, para postular
%    johan-rodriguez-cv-es.tex         → 2 páginas, versión detallada
%
%  Su gemelo en inglés es cv-body.tex. Al cambiar contenido aquí, hay que
%  cambiarlo allí: no se traducen solos.
% =====================================================================
% =====================================================================
%  Preámbulo compartido de los CV de Johan David Rodriguez Castro.
%
%  No se compila solo. Lo cargan los cuerpos: cv-body.tex (EN general),
%  cv-body-es.tex (ES general) y cv-holcim-rpa.tex (ES, dirigido).
%
%  Dos parámetros, que fija el wrapper ANTES de hacer \input de su cuerpo:
%    \cvvariant  = onepage | full   → tamaño de letra y márgenes
%    \cvlang     = en | es          → babel, para partir palabras bien
%
%  Reglas de formato, tomadas de la práctica estándar para perfiles sin
%  experiencia laboral larga y verificadas con `pdftotext` sobre el PDF:
%    · Una sola columna. Nada de cajas de texto ni columnas laterales.
%    · Ninguna tabla de dos columnas para las skills: al extraer el texto,
%      un ATS lee todas las etiquetas primero y después todos los valores,
%      y la asociación etiqueta→keywords se pierde. Se usan líneas simples.
%    · Todo es texto real. Nada dentro de una imagen.
%    · Los proyectos se escriben con la misma estructura que un empleo:
%      título, entidad, fechas y bullets con cifras.
% =====================================================================
\providecommand{\cvvariant}{full}
\def\cvonepagename{onepage}
\ifx\cvvariant\cvonepagename \def\cvfontsize{9pt}\else \def\cvfontsize{10pt}\fi

% extarticle (paquete extsizes) admite 9pt; la clase article estándar no.
\documentclass[\cvfontsize,a4paper]{extarticle}

\newif\ifonepage
\ifx\cvvariant\cvonepagename \onepagetrue \else \onepagefalse \fi

\usepackage[T1]{fontenc}
\usepackage[utf8]{inputenc}
\usepackage{tgtermes}                 % cuerpo: compatible con Times
\usepackage{tgheros}                  % títulos: compatible con Helvetica
\usepackage{microtype}

% Partición de palabras correcta por idioma. es-noshorthands evita que babel
% se apropie de caracteres como " o ., que aquí se usan literalmente.
\providecommand{\cvlang}{en}
\def\cvlanges{es}
\ifx\cvlang\cvlanges
  \usepackage[spanish,es-noshorthands]{babel}
\else
  \usepackage[english]{babel}
\fi


\ifonepage
  \usepackage[a4paper, top=0.85cm, bottom=0.85cm, left=1.2cm, right=1.2cm]{geometry}
\else
  \usepackage[a4paper, top=1.2cm, bottom=1.2cm, left=1.4cm, right=1.4cm]{geometry}
\fi

\usepackage{xcolor}
\definecolor{accent}{HTML}{1F4E79}
\definecolor{subtle}{HTML}{555555}

\usepackage[
  colorlinks=true, urlcolor=accent, linkcolor=accent,
  pdftitle={Johan David Rodriguez Castro — CV},
  pdfauthor={Johan David Rodriguez Castro},
  pdfsubject={Security \& AI Engineer},
  pdfkeywords={security engineering, detection engineering, supply-chain security,
               vulnerability management, applied machine learning, Python, AWS, Terraform}
]{hyperref}

\usepackage{titlesec}
\usepackage{enumitem}
\usepackage{ragged2e}

\titleformat{\section}[block]
  {\sffamily\bfseries\color{accent}\fontsize{10.5}{12}\selectfont}
  {}{0em}{\MakeUppercase}[\vspace{-0.9ex}{\color{accent}\rule{\linewidth}{0.7pt}}]
\titlespacing*{\section}{0pt}{0.9ex plus 0.2ex}{0.6ex}

\setlist[itemize]{
  leftmargin=1.1em, topsep=0.15ex, itemsep=0.05ex, parsep=0pt,
  label=\textcolor{accent}{\textbullet}
}

\setlength{\parindent}{0pt}
\pagestyle{empty}

% ── Entrada de empleo o proyecto ─────────────────────────────────────
% \entry{título}{subtítulo}{fechas · entidad}
% Se usa \hfill dentro de un solo párrafo (no una tabla) para que la
% extracción de texto devuelva título y fecha en la MISMA línea.
\newcommand{\entry}[3]{%
  \vspace{0.5ex}\par\noindent
  \textbf{#1}\ifx\relax#2\relax\else{\,---\,\itshape #2}\fi
  \hfill{\footnotesize\color{subtle}#3}\par\nobreak\vspace{-0.1ex}%
}

% Línea de enlaces bajo una entrada.
\newcommand{\evidence}[1]{{\footnotesize\color{subtle}#1}\par\nobreak}

% Entrada con evidencia: \project{título}{subtítulo}{enlaces}{fechas · entidad}
% En 1 página los enlaces se pliegan en la línea del título para no gastar
% una línea por proyecto.
\newcommand{\project}[4]{%
  \ifonepage
    \entry{#1}{#2}{#3 \ {\color{subtle}$\vert$} \ #4}%
  \else
    \entry{#1}{#2}{#4}\evidence{#3}%
  \fi
}

% Fila de skills: una línea de texto corrido, no una celda de tabla.
\newcommand{\skill}[2]{%
  \par\vspace{0.35ex}\noindent\textbf{#1:}\ #2\par
}

\AtBeginDocument{\setlength{\parindent}{0pt}}

\begin{document}


% =====================================================================
\begin{center}
  {\sffamily\bfseries\color{accent}\fontsize{18}{21}\selectfont
   JOHAN DAVID RODRIGUEZ CASTRO}\\[0.5ex]
  {\footnotesize\color{subtle}
   Bucaramanga, Santander, Colombia
   \ $\vert$ \ +57 313 878 4948
   \ $\vert$ \ \href{mailto:johan.rc2020@gmail.com}{johan.rc2020@gmail.com}}\\[0.4ex]
  {\footnotesize
   \href{https://github.com/Yoyagm}{github.com/Yoyagm}
   \ {\color{subtle}$\vert$} \
   \href{https://www.linkedin.com/in/johan-rodriguez-97bb29323/}{linkedin.com/in/johan-rodriguez-97bb29323}
   \ {\color{subtle}$\vert$} \
   \href{https://johanrodriguez.is-a.dev/es}{johanrodriguez.is-a.dev}}
\end{center}
\vspace{0.2ex}

% =====================================================================
\section{Perfil Profesional}
\justifying
\ifonepage
Estudiante de último semestre de Ingeniería de Sistemas y Software (UPB Bucaramanga), en la
intersección entre ingeniería de seguridad y machine learning aplicado. Entrego herramientas, no
informes: un escáner de cadena de suministro de código abierto, un clasificador de phishing que
auditó su propio benchmark y publicó la cifra honesta en lugar de la que favorecía, y un dashboard
de priorización de vulnerabilidades que desmontó midiendo la ventaja que venía a demostrar.
Inglés C2 verificado (EF SET 80/100).
\else
Estudiante de último semestre de Ingeniería de Sistemas y Software (UPB Bucaramanga), en la
intersección entre ingeniería de seguridad y machine learning aplicado. Entrego herramientas, no
informes: un escáner de cadena de suministro de código abierto, un clasificador de phishing que
auditó su propio benchmark y publicó la cifra honesta en lugar de la que favorecía, y un dashboard
de priorización de vulnerabilidades que desmontó midiendo la ventaja que venía a demostrar. Me
muevo con igual soltura en infraestructura cloud —AWS, Terraform, CI/CD— y en desarrollo
full-stack, con experiencia liderando bajo Scrum e inglés C2 verificado (EF SET 80/100). Busco un
rol de Seguridad, Detection Engineering o Datos e IA.
\fi

% =====================================================================
\section{Competencias Técnicas}
\vspace{0.2ex}
\skill{Seguridad y Detección}{Seguridad de la cadena de suministro de software (slopsquatting),
  detection-as-code, mapeo a MITRE ATT\&CK, priorización de vulnerabilidades basada en riesgo
  (CISA KEV / EPSS / CVSS), análisis de phishing y URLs, 2FA TOTP, SDLC seguro y gates de
  seguridad en CI}
\skill{Datos, ML e IA}{LightGBM, XGBoost, scikit-learn; detección de fugas de datos y diseño de
  evaluaciones honestas; explicabilidad con SHAP; validación agrupada, temporal y walk-forward;
  calibración de probabilidades; NLP con FinBERT; orquestación multiagente con CrewAI; RAG con
  ChromaDB; pandas / NumPy}
\skill{Cloud y DevOps}{AWS (IAM, VPC, EC2, Auto Scaling, ALB, RDS Multi-AZ, ECR, CloudWatch),
  Terraform (IaC), Docker, CI/CD con GitHub Actions, pruebas de carga con k6, Google Cloud
  (BigQuery ML, Cloud Storage, DLP)}
\skill{Programación}{Python, TypeScript / JavaScript, C\#, SQL, VBA, Dart, HCL, LaTeX}
\skill{Bases de Datos}{PostgreSQL, MySQL, SQLite, DuckDB, ChromaDB; arquitecturas híbridas
  SQL / NoSQL}
\skill{Métodos y Herramientas}{Ágil / Scrum, desarrollo dirigido por especificación, pytest,
  Playwright, mypy strict, modelado UML, Git}

% =====================================================================
\section{Experiencia}

\project{Desarrollador / Soporte Técnico}{}%
  {GOATGuard: \href{https://github.com/Yoyagm/goatguard-app}{github.com/Yoyagm/goatguard-app}}%
  {Ene 2026 -- Actualidad \ $\vert$ \ Pontificia Universidad Bolivariana}
\begin{itemize}
  \item Diseñé y construí GOATGuard, una plataforma full-stack de monitoreo de seguridad de red
        (FastAPI + PostgreSQL + Flutter), para cerrar la brecha de movilidad de los administradores
        que gestionan activos tecnológicos distribuidos.
  \item Construí un motor de clasificación automática de activos de red (routers, impresoras,
        cámaras) con monitoreo en tiempo real y dashboards de KPIs, y detección proactiva de port
        scans e intentos de intrusión, aplicando principios de Blue Team y detection engineering.
  \ifonepage\else
  \item Endurecí el cliente móvil con un flujo completo de 2FA TOTP (enrolamiento por QR, códigos
        de respaldo y de recuperación) y JWT en el almacenamiento seguro del sistema operativo,
        tras un interceptor global de errores 401.
  \item Brindo soporte técnico a los usuarios finales de la institución.
  \fi
\end{itemize}

\entry{Scrum Master}{}{Jun 2025 -- Nov 2025 \ $\vert$ \ Pontificia Universidad Bolivariana}
\begin{itemize}
  \item Gestioné el ciclo de vida de desarrollo bajo Scrum, facilitando las ceremonias y
        optimizando el Product Backlog para asegurar entrega incremental de valor.
  \item Contribuí al diseño de una infraestructura de alta disponibilidad orientada a 100.000
        usuarios concurrentes con objetivo de latencia por debajo de 500\,ms, sobre arquitecturas
        de datos híbridas SQL / NoSQL / en memoria.
\end{itemize}

\ifonepage\else
\entry{Auxiliar de Eventos Deportivos}{}{May 2025 -- Actualidad \ $\vert$ \ EM Torneos y Competencias}
\begin{itemize}
  \item Apoyo la organización y ejecución de torneos deportivos y eventos competitivos.
\end{itemize}
\fi

% =====================================================================
\section{Proyectos Técnicos Destacados}

\project{SlopGuard}{Escáner de Cadena de Suministro de Código Abierto}%
  {\href{https://github.com/Yoyagm/slopguard}{github.com/Yoyagm/slopguard}}%
  {2026 \ $\vert$ \ Independiente}
\begin{itemize}
  \item Construí un guardián de pre-instalación que puntúa dependencias de PyPI y npm frente a
        nombres alucinados, typosquatted y maliciosos \emph{antes} de que se ejecute código alguno,
        con un motor de detección de cinco capas y cero dependencias en tiempo de ejecución.
  \ifonepage\else
  \item Diseñé un invariante anti-falso-positivo demostrable: la capa opcional de alucinación por
        LLM es estructuralmente incapaz de bloquear un paquete legítimo por sí sola.
  \fi
  \item Aseguré la corrección con un gate de CI al $\geq$90\,\% de cobertura global ($\geq$95\,\%
        en rutas críticas) y ocho contratos de arquitectura con import-linter, y entregué cuatro
        puntos de integración desde un solo núcleo: CLI, hook de pre-commit, GitHub Action y SaaS
        auto-hospedable.
\end{itemize}

\ifonepage
\project{Detector de Phishing por URL}{Detección Explicable y Auditoría de Fugas}%
  {\href{https://yoyagm.github.io/phishing-url-detector/}{yoyagm.github.io/phishing-url-detector}}%
  {2026 \ $\vert$ \ Independiente}
\else
\project{Detector de Phishing por URL}{Detección Explicable y Auditoría de Fugas del Dataset}%
  {\href{https://yoyagm.github.io/phishing-url-detector/}{yoyagm.github.io/phishing-url-detector}
   \ $\vert$ \ \href{https://github.com/Yoyagm/phishing-url-detector}{github.com/Yoyagm/phishing-url-detector}}%
  {2026 \ $\vert$ \ Independiente}
\fi
\begin{itemize}
  \item Audité PhiUSIIL (UCI id 967, 235.795 URLs) antes de entrenar y medí dos fugas de etiqueta
        independientes: una regla de una línea alcanza 99,67\,\% de exactitud, y features léxicas
        hechas a mano sobre la URL cruda —el arreglo obvio— siguen dando 99,51\,\% sin medir nada.
  \item Reentrené solo sobre nombres de host canonicalizados y publiqué el resultado honesto:
        ROC-AUC 0,937 y 70,5\,\% de recall con 1\,\% de falsos positivos, contra phishing enviado
        dos años después de recolectar los datos de entrenamiento.
  \ifonepage
  \item Lo entregué: TreeSHAP exacto reimplementado en JavaScript sin dependencias (paridad
        verificada a 1e-12 contra la referencia de Python), 113 tests incluyendo guardianes de
        fugas con controles negativos, y un servicio FastAPI containerizado.
  \else
  \item Reimplementé TreeSHAP exacto dependiente del camino en JavaScript sin dependencias, con
        paridad verificada a 1e-12 contra la referencia de Python sobre 28.917 atribuciones, para
        que la demo explique cada veredicto sin backend.
  \item Entregué 113 tests —incluidos guardianes de fugas que vienen con controles negativos de
        fuga plantada—, un servicio FastAPI containerizado y un informe técnico de nueve páginas
        en LaTeX.
  \fi
\end{itemize}

\project{Dashboard de Priorización de Vulnerabilidades}{CISA KEV + EPSS + NVD}%
  {\href{https://github.com/Yoyagm/vuln-prioritization-dashboard}{github.com/Yoyagm/vuln-prioritization-dashboard}}%
  {2026 \ $\vert$ \ Independiente}
\begin{itemize}
  \item Construí una tubería de ingesta reproducible sobre tres feeds públicos que ordena 359.399
        CVEs vivos, con contratos de datos versionados y tests de esquema que rompen la build
        cuando un feed cambia de forma.
  \ifonepage
  \item Reproduje la afirmación de la industria de que EPSS es 6,1$\times$ más eficiente que CVSS,
        y luego demostré con puntuaciones congeladas en cinco cortes que la mayor parte es
        artefacto de evaluación: EPSS gana la cabeza de la cola (1,2--1,8$\times$) y pierde la cola
        larga (0,5--0,7$\times$).
  \else
  \item Reproduje la afirmación de la industria de que EPSS es 6,1$\times$ más eficiente que CVSS,
        y luego demostré —congelando las puntuaciones en cinco fechas pasadas y contando solo la
        explotación confirmada después— que la mayor parte es artefacto de evaluación: EPSS gana la
        cabeza de la cola (1,2--1,8$\times$) y pierde la cola larga (0,5--0,7$\times$).
  \item Calculé intervalos de confianza del 95\,\% por bootstrap sobre los positivos y protegí el
        promediado de empates de la métrica central con un test que viene con control negativo;
        76 tests y cinco ADRs.
  \fi
\end{itemize}

\ifonepage
\vspace{0.6ex}
{\footnotesize\noindent\textbf{Además:} infraestructura AWS de tres capas en alta disponibilidad con
Terraform + GitHub Actions
(\href{https://github.com/Yoyagm/eval-final-dgiti}{eval-final-dgiti}) \ $\vert$ \ tubería autónoma
de ML y agentes para mercados, 505 tests \ $\vert$ \ GreenPulse, plataforma agronómica offline-first.
Casos de estudio completos:
\href{https://johanrodriguez.is-a.dev/es/proyectos}{johanrodriguez.is-a.dev/proyectos}.\par}
\else

\project{Infraestructura AWS de Alta Disponibilidad}{Terraform, FastAPI y CI/CD}%
  {\href{https://github.com/Yoyagm/eval-final-dgiti}{github.com/Yoyagm/eval-final-dgiti}}%
  {2026 \ $\vert$ \ UPB}
\begin{itemize}
  \item Codifiqué íntegramente en Terraform una arquitectura de tres capas repartida en dos zonas
        de disponibilidad: VPC, ALB, Auto Scaling Group, RDS PostgreSQL Multi-AZ, ECR y CloudWatch.
  \item Apliqué mínimo privilegio con una cadena estricta de security groups: la base de datos no
        tiene acceso público ni egress, y solo es alcanzable desde la capa de aplicación.
  \item Automaticé el despliegue con GitHub Actions (linters, tests, build, push a ECR e instance
        refresh del ASG), manteniendo Terraform fuera de la ruta caliente, y validé el balanceo
        bajo carga con k6 contra métricas de CloudWatch.
\end{itemize}

\project{Sistema Autónomo de Inteligencia de Mercado}{ML Aplicado e IA Multiagente}%
  {Repositorio privado --- recorrido disponible a petición}%
  {2025 -- Actualidad \ $\vert$ \ Independiente}
\begin{itemize}
  \item Construí una tubería dirigida por eventos que integra APIs de datos de mercado, cálculo de
        indicadores técnicos y un modelo de sentimiento FinBERT, orquestando cuatro agentes CrewAI
        en informes automáticos diarios.
  \item Diseñé una capa de gobernanza con siete comprobaciones de riesgo más un cortacircuitos por
        drawdown de cartera; las señales de confianza media escalan a aprobación humana por
        Telegram y expiran sin respuesta.
  \item Entrené un ensemble XGBoost / RandomForest / LogisticRegression con validación
        walk-forward, hueco purgado de 21 días y calibración de Platt; validado por 505 tests
        automatizados en 17 suites.
  \item Opera en paper trading de Alpaca y se queda ahí deliberadamente: una puerta GO/NO-GO
        bloquea el capital real hasta que el modelo supere el 55\,\% de exactitud (hoy 52,9\,\%).
\end{itemize}

\entry{GreenPulse}{Plataforma Agronómica de Monitoreo Offline-First}{2025 \ $\vert$ \ Proyecto de Grado}
\begin{itemize}
  \item Migré y rediseñé el backend a FastAPI con una capa de sincronización offline-first sobre
        SQLite y un catálogo de 35 especies con escaneo por código QR; documentado en formato IEEE.
\end{itemize}
\fi

% =====================================================================
\section{Formación}

\entry{Ingeniería de Sistemas y Software}{}{Feb 2023 -- Actualidad \ $\vert$ \ Pontificia Universidad Bolivariana}
\begin{itemize}
  \ifonepage
  \item Último semestre. Cursos avanzados de arquitectura de software y diseño de sistemas; labs
        1--6 de AWS Academy Cloud Foundations (IAM, VPC, EC2, EBS, Auto Scaling, balanceo).
  \else
  \item Último semestre. Cursos avanzados de arquitectura de software, diseño de sistemas y buenas
        prácticas de la industria.
  \item Formación AWS Academy Cloud Foundations --- labs 1--6 completados: IAM, VPC, EC2, EBS,
        Auto Scaling y balanceo de carga, con implementaciones de IaC en Terraform (2026).
  \fi
\end{itemize}

% =====================================================================
\section{Certificaciones}

\ifonepage
{\footnotesize\noindent
\textbf{Google Cloud} (2026) --- 7 skill badges: Secure Lakehouse Data \textbullet\ Protect Sensitive
Data with Data Loss Prevention \textbullet\ Create ML Models with BigQuery ML \textbullet\ Use APIs to
Work with Cloud Storage \textbullet\ Gemini Enterprise (flujos multiagente, primera aplicación)
\textbullet\ Vibe Coding and MCP. \ \textbf{Google} (2026) --- DeepMind: Train A Small Language Model.\par
\vspace{0.2ex}\noindent
\textbf{Cisco Networking Academy} (2026) --- Introduction to Cybersecurity. \
\textbf{Hugging Face} (2026) --- AI Agents Course. \
\textbf{freeCodeCamp} (2026) --- Machine Learning with Python \textbullet\ Data Analysis with Python
\textbullet\ Scientific Computing with Python \textbullet\ Data Visualization \textbullet\ JavaScript
Algorithms and Data Structures.\par
\vspace{0.2ex}\noindent
\textbf{Idiomas:} Español --- Nativo \textbullet\ Inglés --- C2 Competente
(\href{https://cert.efset.org/es/qpqdd3}{EF SET 80/100, verificado}).
{\color{subtle}Todas verificables públicamente en
\href{https://www.linkedin.com/in/johan-rodriguez-97bb29323/details/certifications/}{linkedin.com/in/johan-rodriguez-97bb29323}.}\par}
\else
{\footnotesize\color{subtle}Todas las credenciales son verificables públicamente; lista completa con
enlaces en
\href{https://www.linkedin.com/in/johan-rodriguez-97bb29323/details/certifications/}{linkedin.com/in/johan-rodriguez-97bb29323}.}
\vspace{0.4ex}
\begin{itemize}
  \item \textbf{Google Cloud} (2026) --- Secure Lakehouse Data \textbullet\ Protect Sensitive Data
        with Data Loss Prevention \textbullet\ Create ML Models with BigQuery ML \textbullet\ Use
        APIs to Work with Cloud Storage \textbullet\ Orchestrate Multi-agent Workflows with Gemini
        Enterprise \textbullet\ Create Your First Gemini Enterprise Application \textbullet\ Build a
        Smart Cloud Application with Vibe Coding and MCP \textbullet\ Google DeepMind: Train A Small
        Language Model
  \item \textbf{Cisco Networking Academy} (2026) --- Introduction to Cybersecurity
  \item \textbf{Hugging Face} (2026) --- AI Agents Course
  \item \textbf{freeCodeCamp} (2026) --- Machine Learning with Python \textbullet\ Data Analysis
        with Python \textbullet\ Scientific Computing with Python \textbullet\ Data Visualization
        \textbullet\ JavaScript Algorithms and Data Structures
  \item \textbf{EF SET} (2026) --- Certificado de inglés, MCER C2 (Competente), 80/100 global
        (lectura 84, comprensión oral 89)
\end{itemize}
\fi

% =====================================================================
\ifonepage\else
\section{Idiomas}
Español --- Nativo \ \ {\color{subtle}$\vert$} \ \ Inglés --- C2 Competente
(\href{https://cert.efset.org/es/qpqdd3}{EF SET 80/100, verificado})
\fi

\end{document}
